\documentclass[11pt,letterpaper]{article}
\usepackage[T1]{fontenc}
\usepackage{lmodern,amsmath,amssymb,amsthm,microtype}
\usepackage[margin=1in]{geometry}
\usepackage[hidelinks]{hyperref}
\newtheorem{theorem}{Theorem}
\title{No mixed unit edges after adjoining a product rotation}
\author{Aafi Khan Malek\\\small Independent Researcher, Ahmedabad, India\\
\small\href{mailto:aafimalek2032@gmail.com}{aafimalek2032@gmail.com}}
\date{November 2026}
\begin{document}
\maketitle
Put \(s=\sqrt{-3}\), \(r=\sqrt{33}\), \(h=\sqrt{-7}\), using
the positive real or positive imaginary root. Let
\[
 \omega=\frac{1+s}{2},\quad \Lambda=\mathbb Z[\omega],\quad
 \alpha=\frac{r+s}{6},\quad \beta=\frac{3+h}{4}.
\]
Both complex rotations have norm one. Our
\href{https://github.com/Aafimalek/novel-discoveries/tree/8be0f5af188954944d37043740b20aeaf1a0c2bf/hadwiger-nelson/verification/parts-native-field-obstruction/mixed-additive-lattice}{earlier three-factor calculation}
classified unit differences in \(\Lambda+\alpha\Lambda+\beta\Lambda\).
The product rotation \(\alpha\beta\) lies outside that additive
group. Even adjoining all of \(\alpha\beta\Lambda\) does not
create any edges combining different factors.

\begin{theorem}
The unit-distance graph on
\[
 M_4=\Lambda+\alpha\Lambda+\beta\Lambda+\alpha\beta\Lambda
\]
is the Cartesian product of four triangular lattices. Its
chromatic number is three. Any chosen distinct nonadjacent pair
can be made monochromatic in a proper four-coloring of \(M_4\).
\end{theorem}
\begin{proof}
The square classes of \(-3,33,-7\) are independent: the primes
11 and 7 detect the last two classes, and \(-3\) is not a rational
square. In particular \(1,r,h,rh\) are independent over
\(\mathbb Q(s)\). For \(t,u,v,w\in\Lambda\), write
\[
\begin{split}
 d &=t+\alpha u+\beta v+\alpha\beta w,\qquad
 24d=a+rb+hc+rh e,\\
 a&=24t+4su+18v+3sw,\quad b=4u+3w,\quad
 c=6v+sw,\quad e=w.
\end{split}
\]
Independence followed by these triangular coefficient equations
gives uniqueness of the four-factor representation.

Write \(N(z)=z\overline z\). Conjugation fixes \(r\) and
negates \(h\); the rational coefficient of \(N(24d)\) is
\[
 N(a)+33N(b)+7N(c)+231N(e).
\]
The terms in the independent directions \(r,hs,rhs\) vanish
when the norm is rational, giving
\[
\begin{split}
 \Re(a\overline b)+7\Re(c\overline e)&=0,\\
 \Im(c\overline a)+33\Im(e\overline b)&=0,\\
 \Im(e\overline a)+\Im(c\overline b)&=0.
\end{split}
\]
For a unit displacement the rational coefficient is 576. If
\(w=0\), the earlier three-factor theorem shows the edge changes
just one factor. Otherwise \(231N(w)\le576\), so \(N(w)=1\)
or 2. The latter is impossible:
\(N(m+n\omega)=m^2+mn+n^2\equiv(m-n)^2\pmod3\).
Multiplication of all summands by \(\overline w\) preserves their
lattices and normalizes \(e=w=1\).

Then \(b\equiv3\pmod{4\Lambda}\) and
\(N(b)\le345/33<11\). A nonzero \(\omega\) coordinate
in \(b\) is a nonzero multiple of four, with imaginary part
of squared magnitude at least 12. Thus \(b\in\mathbb Z\),
and the bound leaves only \(b=3\) or \(b=-1\).

If \(b=3\), then \(u=0\) and \(N(c)\le48/7\). Since
\(c=6v+s\), any nonzero \(v\in\Lambda\) would give
\(|c|\ge 6-\sqrt3>\sqrt{48/7}\). Hence \(v=0\) and
\(c=s\). The norm equation gives \(N(a)=27\); with
\(a=24t+3s\), the inequality \(24-3\sqrt3>3\sqrt3\)
implies \(t=0\). The step is entirely in the
\(\alpha\beta\) factor.

If \(b=-1\), then \(u=-1\). The three displayed conditions
reduce to \(\Re(a)=7\Re(c)\), \(\Im(c\overline a)=0\),
and \(\Im(a)=-\Im(c)\). For \(c=x+iy=6v+s\), the number
\(y=(3n+1)\sqrt3\) for some \(n\in\mathbb Z\) is nonzero.
With \(a=7x-iy\), the middle condition says \(8xy=0\).
Thus \(x=0\) and \(N(a)=N(c)\). Now \(v\in\Lambda\) is
purely imaginary, so \(v=ks\) with \(k\in\mathbb Z\).
The norm identity gives \(8N(c)=576-231-33=312\), and
\(c=(6k+1)s\) would require \((6k+1)^2=13\), impossible.

In all cases a unit edge changes exactly one factor by a
lattice unit. Conversely any such change has norm one. This
proves the product description. Color each factor's site
\(m+n\omega\) by \(m+2n\) modulo three, and use the sum
modulo three on the product. A triangle in one factor shows
three is sharp. For any one specified nonadjacent pair, give
its two endpoints the same fresh fourth color.
\end{proof}

\paragraph{Scope.}
No subset of this additive group can supply a six-chromatic
unit-distance witness or force a chosen nonadjacent pair apart
in five colors. This does not cover the ring generated by
\(\alpha,\beta\); for instance
\(\alpha^2=1+(s/3)\alpha\) and
\(\beta^2=-1+(3/2)\beta\) lie outside this integral additive
module. It also does not improve the chromatic bounds for the
full plane. This written proof was developed with AI assistance
and was not formalized in Lean.
\end{document}
